% 原文 p.321
\hypertarget{srcpage:321}{}\label{srcpage:321}
耐故障性をもつ量子計算の詳細には立ち入らないが（この主題のわかりやすい総説として、たとえば文献\cite{ref056,ref367}参照）、しきい値定理とよく呼ばれる重要な結果\cite{ref360,ref368,ref369,ref370,ref371}を述べよう。耐故障性をもつ量子誤り訂正を実現すれば、ある物理的実現におけるゲート操作当たりの誤り率が、あるしきい値より小さいかぎり、\footnote{この精度のしきい値の現在の見積もりは、計算手順当たり約$10^{-2}$から$10^{-6}$まで幅がある。異なる見積もりの分析については文献\cite{ref372}参照。}少なくとも原理的には、任意の大きさの量子計算を効率よく実行できるはずだと示せる。この重要な定理は、誤りの性質について、少数の物理的に妥当な仮定を置けば証明でき、量子コンピューターの実現に基本的な障害はないはずだということを示唆する。

\subsectionlink{デコヒーレンスのない部分空間での量子計算}
デコヒーレンスのない部分空間（DFS）\cite{ref089,ref090,ref091,ref092,ref093,ref094,ref095,ref096,ref097,ref098}、またはポインター部分空間\cite{ref009}の一般的概念は、2.8.1節で導入した。DFSとは、系のHilbert空間$\mathcal H$の部分空間$\mathcal H_{\mathrm{DF}}$であり、その中の任意の状態$|\psi\rangle$が、環境との相互作用にもかかわらず、（厳密に、または近似的に）純粋なまま保たれるものであった。系を計算用量子ビットの集まりとする量子計算に適用すると、DFSにより、Hilbert空間の「静かな片隅」に量子情報を符号化して、環境の有害な影響から、その情報を自動的に保護できる。量子誤り訂正と異なり、DFSは誤りがそもそも起こることを防ぐため、本質的な誤り回避の方策を表す。

相互作用Hamiltonianの一般的な対角分解
\begin{equation}\hat H_{\mathrm{int}}=\sum_\alpha\hat S_\alpha\otimes\hat E_\alpha\tag{7.48}\end{equation}
が与えられるとき、DFSが現れるための必要条件は、演算子$\hat S_\alpha$の正規直交する縮退固有状態$|s_i\rangle$の集合が存在すること、すなわち
\begin{equation}\hat S_\alpha|s_i\rangle=\lambda^{(\alpha)}|s_i\rangle\quad\text{すべての$\alpha$と$i$について}\tag{7.49}\end{equation}
であることを示した［式(2.96)参照］。2.8.1節では、この条件が力学的対称性の存在に相当することも論じた。以下では、すべての量子ビットが、まったく同じ仕方で一つの環境と結合する模型を考える。
% 原文 p.322
\hypertarget{srcpage:322}{}\label{srcpage:322}
この場合、系と環境の相互作用は、量子ビットの任意の置換に対して完全に対称であり、最大の大きさのDFSをもたらす。独立な小さい結合項が追加されて、この完全な対称性が破れると、何が起こるだろうか。DFSには、このような摂動に対して、十分な頑健さがあってほしい。Lidar、Chuang、Whaley\cite{ref090}による初期の研究に続き、Bacon、Lidar、Whaley\cite{ref099}は、摂動の強さの一次まででは、DFS内の量子情報の保存は、このような摂動に対して、時間についてのすべての次数で安定であることを示した。しかし、DFSに符号化された量子情報の処理は、対称性を破る摂動に対して、時間の一次までしか頑健でないことがわかった。著者たちは、このDFS上の量子計算の不十分な頑健さを、DFSと能動的な量子誤り訂正方式を組み合わせることで、どのように補えるかも示した\cite{ref094}。

\subsubsectionlink{デコヒーレンスのない部分空間はどのようなものか}
これまではDFSの存在の一般条件を述べるだけで、その構造の具体例は与えなかった。また、DFSが小さすぎて、必要な高次元の重ね合わせ状態を実効的に作れないのであれば、量子計算にはほとんど役立たないことは明らかである。したがって、系の全Hilbert空間の次元に対して、DFSがどの程度の大きさをもつかという問題は非常に重要である。意外ではないだろうが、この問いに具体的に答えるには、系と環境の相互作用について特定の模型を採用しなければならない。ここでは、DFSの文献（たとえば文献\cite{ref089,ref091,ref093,ref096,ref237}参照）で最も一般的に考えられる模型、すなわち5.3節で論じたスピン・ボソン模型を考える。

\subsubsectionlink*{独立なデコヒーレンスと集団的デコヒーレンス}
量子ビットのデコヒーレンスを模型化する二つの極限を区別しよう。第一の極限は独立なデコヒーレンスであり、各量子ビットが独立に自分自身の環境と結合し、これらの環境間には相互作用がないという仮定である。この場合は7.4.2節ですでに述べた。独立なデコヒーレンスの仮定は、量子誤り訂正の文脈でしばしば置かれる。それは、量子ビットに作用する誤り過程が完全に無相関であることを意味する。したがって、ある特定の誤りが一つの量子ビットに起こる確率が$p$なら、この誤りが$K$個の量子ビットで起こる確率は$p^K$である。この場合、量子ビット誤りは、単に重み1の誤り演算子の線形結合、すなわち、他のすべての量子ビットは影響を受けない単一量子ビット誤りの組合せとして表せる。多くの誤り訂正方式（7.4.5節で論じた三量子ビット符号など）は、このような単一量子ビット誤りの訂正でのみ効率がよいため、それらの方式の基礎には、しばしば独立なデコヒーレンスの仮定がある。

しかし、量子ビットが（環境の典型的なコヒーレンス長と比べて）互いに空間的に近い場合、この仮定はあまり現実的でない。
% 原文 p.323
\hypertarget{srcpage:323}{}\label{srcpage:323}
その場合、すべての量子ビットは（近似的に）同じ環境を「感じ」、誤りは複数の量子ビット間で相関しやすい。この状況に対応する極限は集団的デコヒーレンスであり、すべての量子ビットがまったく同じ環境と結合する。

完全に独立なデコヒーレンスと、純粋に集団的なデコヒーレンスという両極端が、DFSの大きさの限界を定めることがわかる。スピン・ボソン模型では、集団的デコヒーレンスの極限で、量子ビットが多い場合、DFSの次元は元のHilbert空間の次元へ漸近的に近付き、最適な、デコヒーレンスのない符号化が可能になる。しかし反対の、独立なデコヒーレンスの極限では、論理量子ビットを符号化できるDFSが存在しないことを見る。これらの結果を確立するための詳細な模型と証明は再現しないが、それでもDFSの大きさと特定のデコヒーレンス模型の関係について、ある程度の直観を得たい。以下ではそれを行う。

\subsubsectionlink*{簡単な例}
まず、スピン・ボソン模型における集団的デコヒーレンスの場合を考えよう\cite{ref089,ref091,ref093,ref096,ref237}。ここではDFSは、ほとんど定義によって現れる。$N$量子ビットの集まりの相互作用Hamiltonianは、式(5.90)の相互作用Hamiltonianの一般化、すなわち
\begin{equation}\hat H_{\mathrm{int}}=\sum_{i=1}^N\hat\sigma_z^{(i)}\otimes\sum_j(g_{ij}\hat a_j^\dagger+g_{ij}^{*}\hat a_j)\equiv\sum_{i=1}^N\hat\sigma_z^{(i)}\otimes\hat E_i\tag{7.50}\end{equation}
で与えられる。ここで$g_{ij}$は、第$i$量子ビットと第$j$環境振動子の結合の強さを表し、$\hat a_j$（$\hat a_j^\dagger$）は、ボソン場の第$j$モードの消滅（生成）演算子である。

集団的デコヒーレンスの仮定は、すべての量子ビットが同じ環境と結合することを意味する。そのため、結合$g_{ij}$（したがって環境演算子$\hat E_i$）は添字$i$に依存してはならない。すると式(7.50)は
\begin{equation}\hat H_{\mathrm{int}}=\left(\sum_i\hat\sigma_z^{(i)}\right)\otimes\hat E\equiv\hat S_z\otimes\hat E\tag{7.51}\end{equation}
となる。環境演算子$\hat E$の正確な形は、ここではそれほど重要でない。主要な点は、相互作用Hamiltonianが式(7.51)の形に書かれ、和が一つの項$\hat S_z\otimes\hat E$に置き換わったことである。

ここで、DFSは、相互作用Hamiltonianの系演算子$\hat S_\alpha$の、縮退した固有状態の集合によって張られることを思い出そう［式(7.49)参照］。したがって、この場合のDFSは、式(7.51)で定義した集団スピン演算子$\hat S_z$の縮退固有状態によって張られる。
% 原文 p.324
\hypertarget{srcpage:324}{}\label{srcpage:324}
計算基底(7.5)で表された任意の$N$量子ビット積状態は、明らかに$\hat S_z$の固有状態である。可能な整数固有値$m$には、$m=-N$（基底状態$|1\cdots1\rangle$に対応）から$m=+N$（基底状態$|0\cdots0\rangle$に対応）まで、$2N+1$個の異なる値がある。一般に、$m_0$個の量子ビットが「0」状態にある計算基底状態の固有値は、$m=2m_0-N$である。

$\hat S_z$の同じ固有値$m$をもつ、互いに直交する計算基底状態の個数が最大になるのは、$m=0$の基底状態の集合$S_0$、すなわち「0」状態と「1」状態の量子ビットが等しい個数含まれる状態の場合である。形式的には、
\begin{equation}S_0=\left\{|i_1\cdots i_N\rangle\ \middle|\ i_j\in\{0,1\},\ \sum_{j=1}^N i_j=0\right\}\tag{7.52}\end{equation}
と書け、この集合の状態数は
\begin{equation}n_0=\binom{N}{N/2}\tag{7.53}\end{equation}
である。したがって、$S_0$に含まれる状態は、次元$n_0$のDFSを張る。$N$が大きいとき、Stirlingの公式を用い、式(7.53)の二項係数を
\begin{equation}\log_2\binom{N}{N/2}\approx N-\frac12\log_2(\pi N/2)\xrightarrow{N\gg1}N\tag{7.54}\end{equation}
と近似できる。したがって、先に述べたとおり、集団的デコヒーレンスの極限で、DFSの次元は元のHilbert空間の次元に近付き、符号化効率は1に近付く（すなわち、$N\longrightarrow\infty$の極限では、元のHilbert空間の$2^N$個すべての量子ビット基底状態を、デコヒーレンスなしに符号化できる）。

たとえば$N=4$量子ビットでは、計算基底状態の集合
\begin{equation}S_0=\{|0011\rangle,|0101\rangle,|0110\rangle,|1001\rangle,|1010\rangle,|1100\rangle\}\tag{7.55}\end{equation}
が、次元6の、最大の大きさのDFSを張る。これに対し、元のHilbert空間の次元は$2^4=16$である。したがって、ここで考えた集団的デコヒーレンス模型では、DFSの中に二つまでの論理量子ビットを符号化できる。三つの論理量子ビットには、すでに$2^3$個の計算基底状態が必要であり、式(7.55)は明らかにそれを与えない。そのため、四つの物理量子ビットを用いれば、二つ以下の論理量子ビットを必要とする、デコヒーレンスのない量子計算を実行できる。

逆に、純粋に独立なデコヒーレンスの場合には、次元が1より大きいDFSを見つけることはできない。これは次の簡単な観察からわかる。相互作用Hamiltonian(7.50)に現れる環境演算子$\hat E_i$は、今度は一般に、すべて互いに異なる。DFSを見つけるには、通常の方法(7.49)に従い、すべての$i$と$1\leq j\leq N$について
% 原文 p.325
\hypertarget{srcpage:325}{}\label{srcpage:325}
\begin{equation}\left\{\hat I^{(1)}\otimes\cdots\otimes\hat I^{(j-1)}\otimes\hat\sigma_z^{(j)}\otimes\hat I^{(j+1)}\otimes\cdots\otimes\hat I^{(N)}\right\}|s_i\rangle=\lambda^{(j)}|s_i\rangle\tag{7.56}\end{equation}
となる正規直交基底状態$\{|s_i\rangle\}$の集合を求める。読者が容易に確認できるように、この固有値問題を満たす状態は、計算基底状態$|0\cdots0\rangle$という、ただ一つの状態だけである。一つの論理量子ビットを符号化するにも、少なくとも二次元の部分空間が必要なので、スピン・ボソン模型の独立なデコヒーレンスの場合には、量子計算のためのDFSは存在しえない。ポインター部分空間の言葉では、厳密なポインター状態は一つしかなく、系の、この環境によって超選択された優先状態は、基底状態$|0\cdots0\rangle$となる。

もちろん、現実的な状況では、純粋に独立なデコヒーレンスという仮定も、完全に集団的なデコヒーレンスという極限も、完全には適切でない。しかし幸い、この二つの極限それぞれでデコヒーレンスに対抗する、強力な相補的方法がいまや二つある。DFSへの符号化によって、量子ビット系を集団的デコヒーレンス効果から保護し、7.4節で述べた能動的な誤り訂正法によって、独立なデコヒーレンスに起因する単一量子ビット誤りから回復できる。実際、両方法は結び付ける（または「連接する」\cite{ref094}）ことができ、それによって単一量子ビット誤りへの制限を外しても、耐故障性をもつ普遍量子計算が可能になる\cite{ref095,ref373}。

\subsubsectionlink{デコヒーレンスのない部分空間の実験的実現}
DFSの存在を明示的に実証した最初の実験の結果は、Kwiatと共同研究者が2000年に報告した\cite{ref374}。この実験は、量子計算に関係するDFSを実際に実現したというより、原理の実証であったことを強調すべきである。単に、一対の光子のある特定のBell状態［式(2.7)参照］が、集団的デコヒーレンスに対して本質的に免疫をもつ一方、他の三つのBell状態の光子対は、理論予測と一致して、大きなデコヒーレンスを受けることが示された。この実験では、人工環境として、二つの光子の経路に対して調整可能な角度に向けた、小さな二片の石英を用い、その中を通過する光子に、制御できるデコヒーレンス相互作用を生じさせた。二片の石英を同じ角度（$90^\circ$を法として）にそろえると、集団的デコヒーレンスが生じた。角度を変えることは、異なる基底でデコヒーレンスを生じさせることに対応した。本質的には、この実験は単一のポインター状態の存在を示しただけである。

DFSに関する別の実験の結果は、NISTのKielpinskiらが2001年に報告した\cite{ref375}。彼らは、前の段落で述べたKwiatらの実験\cite{ref374}と同様の、二量子ビットBell状態を用いた。しかし今回は、捕獲され、強く相互作用する一対の$^9\mathrm{Be}^+$イオンのDFSへ、一つの論理量子ビットを符号化した。この実験は、符号化された基底状態の重ね合わせが、実際に集団的デコヒーレンスの影響を受けないことを示した。観測されたコヒーレンスの減衰は、符号化されていない量子ビット状態に比べてはるかに遅かった。
% 原文 p.326
\hypertarget{srcpage:326}{}\label{srcpage:326}
その減衰は、集団的デコヒーレンス以外の原因（読み出しパルスの劣化や、イオンの運動状態の加熱など）に帰属できた。したがってこの実験は、物理量子ビットの四次元Hilbert空間を用いて、集団的デコヒーレンスの影響を受けない二次元部分空間を実験的に構築することに成功した。

それ以来、実験は、量子ビットの異なる物理的実現を用い、ますます大きなDFSを作ろうとしてきた。たとえばViolaらは2001年に、NMR量子ビットを用い、初めて三量子ビットDFSの生成に成功した\cite{ref376}。

\subsubsectionlink{環境工学と動的分離}
実用上、すなわち実験的に関心のある状況で、少なくとも近似的にでも、DFSの存在条件をどこまで満たせるかと疑問に思うかもしれない。前節では、実際にDFSを実現した近年の実験をいくつか説明した。しかし、そのDFSはまだ非常に小さい。相当に大きなDFSが存在するには、系と環境の相互作用が、十分に高い対称性を示さなければならないことは明らかであり、与えられた実験状況で、こうした条件が自然に生じる可能性は低いように思われる。

この制限を克服する可能な方法の一つは、環境工学に基づく。ここでは、系と環境の相互作用の構造に、ある種の対称性を能動的に作り出そうとする。明らかな理由から、この方策はときに「環境の対称化」とも呼ばれる。たとえばDalvit、Dziarmaga、Zurek\cite{ref102}は、系と環境の結合を適切に設計して対称化すれば、Bose--Einstein凝縮体の重ね合わせ状態を、相互作用Hamiltonianの（近似的に）縮退した固有状態に対応させられることを示した。したがって、こうした重ね合わせ状態はDFS内に属し、その寿命を大きく延ばすことになる（6.4.1節も参照）。

Poyatos、Cirac、Zoller\cite{ref377}が示したように、捕獲イオンの実効的な系・環境相互作用Hamiltonianのパラメーターを変えることで（制御するレーザー光が環境の役割を果たす）、異なるポインター部分空間を選び、系（捕獲イオン）がどのDFSへ駆動されるかを能動的に制御することもできる。この特徴はその後、デコヒーレンスの制御と模擬への応用として、イオントラップ実験で実証された\cite{ref125,ref353}（文献\cite{ref378}も参照）。これは、ポインター基底という概念に直接の実験的証拠を見事に与える。「古典的」状態の形は、量子力学の外から持ち込まれた、何らかの先験的基準であらかじめ定められているのではなく、2.8節で述べたように、系と環境の相互作用Hamiltonianの構造によって動的に定まる。この相互作用Hamiltonianを変えると、異なる準古典的状態に対応する、さまざまなポインター基底が現れる。

DFSを能動的に作り出す別の方法は、動的分離として知られるようになった\cite{ref379,ref380,ref381,ref382,ref383,ref384}。ここで基本的な工夫は、環境の影響を実効的に打ち消すよう、系のHamiltonianへ時間依存の変更を加えることである。
% 原文 p.327
\hypertarget{srcpage:327}{}\label{srcpage:327}
これらの変更は、系だけに作用する、高速な射影測定列や強い制御場パルス列の形を取る（「量子バンバン制御」\cite{ref379}）。系と環境の相互作用Hamiltonianの構造をまったく知らなくても、分離用の制御場の速度を無限に速くする極限では、デコヒーレンスを任意によく抑制でき、その結果DFSを動的に作り出せる（この場合、それは動的に分離された部分空間を表す）。制御速度が有限である現実的な場合にも、制御速度が、環境もつれの形成率によって定まる最も速い時間尺度を上回れば、この分離技法によって、十分な（完全ではないが）デコヒーレンスからの保護を達成できる。

\subsectionlink{まとめと展望}
量子コンピューターの実現にとって主要な障害は、必要とされる遮蔽と制御可能性のトレードオフである。一方では、デコヒーレンスを最小化するため、望ましくない環境との相互作用から十分に保護された量子ビット系を構築する必要がある。他方では、量子計算の実行に必要なコヒーレント制御を可能にするため、これらの系は十分に開かれたままでなければならない。

このように量子ビットを単に孤立させられない以上、量子計算でデコヒーレンスを最小化するには、能動的な量子誤り訂正、デコヒーレンスのない部分空間、環境工学、動的分離などの技法を、つねに組み合わせる必要がある。単一量子ビット誤りだけに限るという（通常、あまり現実的でない）制限を外すと、耐故障性をもつ普遍量子計算を達成するには、DFSと量子誤り訂正の両方を用いる必要があることを示せる\cite{ref373}。そこでの考えは、まず、動的な対称性を利用する（あるいは作り出す）ことで、十分な程度に近似的なデコヒーレンスのない部分空間へ符号化できるように、量子コンピューターを構築することである。量子ビット系の時間発展に外部から時間依存の変更を加えることで、量子計算の途中にも、実効的なDFSを作り、維持できる。DFSは理想条件からの小さい摂動にある程度頑健なので、この方法は、量子計算にデコヒーレンス効果に対する、ある程度の安定性をもたらす。その後、第二の段階として、量子誤り訂正方式を用い、デコヒーレンスの影響から能動的に回復しようとすることができる。

現在までにコヒーレントに制御できているのは、ほんの一握りの量子ビットであり、有用な量子計算を行うには、まったく足りない。デコヒーレンス時間を十分に延ばし、動作する量子コンピューターの構築を可能にできるか、またそれがいつになるかは、なお今後を待たなければならない。いずれにせよ、量子計算への関心に促された、デコヒーレンスに対抗する方策の研究は、デコヒーレンス自体の理論的・実験的理解を深めるうえでも、実り豊かであった。
% 原文 p.328
\hypertarget{srcpage:328}{}\label{srcpage:328}
デコヒーレンス研究の初期には、量子測定と、一般的な量子から古典への移行に対する基礎的含意に焦点が当てられ、関心の中心は、デコヒーレンスの実際上不可逆な性質と、その極端な効率のよさであった。過去10年のあいだに、量子計算の分野は、デコヒーレンスを生む相互作用の能動的な制御と設計に焦点を当てることで、この範囲を大きく広げた。たとえば、系と環境の相互作用の構造を設計・操作し、デコヒーレンスの影響の小さいHilbert空間の隙間を利用し、作り出す方法を学び、実験的に実証してきた。量子誤り訂正は、デコヒーレンス効果の影響と軽減を、離散的な誤りという言葉で表すようにした。これらの誤りは、デコヒーレンスを引き起こすまさにその効果、すなわち量子もつれを利用することで、実効的に「取り消す」ことができる。さらに多くの例をここに挙げることができ、そのいくつかは本章で論じた。この意味で、デコヒーレンスについての理解は、近年の量子計算の研究から大きな刺激を受けてきた。この研究分野から、デコヒーレンスについての、さらに多くの刺激的な洞察が得られることを期待できる。
